\section{总结与延伸}

本讲从异构部署开始，沿着软件 delivery、runtime substrate、decision chain、制造 workflow 和制度反馈逐层展开。Apollo 把软件与环境状态纳入持续控制，Rubix 把安全和生命周期纳入统一 substrate，Maven 案例说明模型必须进入有权限和反馈的 OODA loop，制造案例说明 workflow 要在约束变化时重算，privacy-security frontier 则要求能力提升与正当程序同时存在。

\begin{importantbox}{九条可迁移结论}
\begin{enumerate}
  \item 描述 scale 时同时统计环境异构性、变更频率和失联时间，不只统计节点数。
  \item 把软件部署写成 desired-state control loop，而不是一串不可观察 shell scripts。
  \item 让 service owner 声明 health、dependency 和 schema compatibility，自动化才能安全。
  \item 用 release channel 与 canary 表达置信度，不追求所有环境同步变更。
  \item Compliance-as-code 应阻止违规状态并生成证据，不只是自动填审计表。
  \item Ephemeral infrastructure 缩短 drift 和持久化窗口，但要求应用支持 failover。
  \item 模型 inference 只是 decision chain 一环，权限、workflow 和 outcome 才决定部署价值。
  \item 现场 spreadsheet 是产品需求和反馈断点的证据，不应简单禁止。
  \item 技术可以扩大 privacy/security 可行集合，合法目标仍需制度选择和监督。
\end{enumerate}
\end{importantbox}

\subsection{实践作业：设计一个异构 Fleet 的安全升级}

假设你维护一个包含公有云、客户机房和断网边缘站点的 30-service 产品，需要修复一个高危 library 漏洞。请设计 catalog metadata、release channels、health monitors、rollback 和例外流程。作业不能只写“使用 Kubernetes”和“自动部署”，必须给出 environment constraints 与数据兼容性。

\begin{enumerate}
  \item 列出至少四类环境，并定义连接性、维护窗口、安全级别和可用容量。
  \item 为受影响 component 建立 dependency、schema、vulnerability 和 owner 记录。
  \item 设计 canary→connected→regulated→air-gapped 的 rollout，写出每阶段进入/退出条件。
  \item 选择一个无法立即升级的环境，给出 compensating control、风险 owner 和到期时间。
  \item 模拟新版本失败，说明如何判断 rollback 安全、怎样处理已经写入的数据。
\end{enumerate}

\begin{knowledgebox}{验收标准}
高质量方案应能回答“现在每个环境运行什么”“为什么它可以升级”“失败后谁被通知”“谁批准例外”“哪条证据证明风险下降”。如果答案依赖某位资深工程师记住所有顺序，说明控制面还没有真正形成。
\end{knowledgebox}

\subsection{拓展阅读}

建议直接阅读本讲目录中的 `source-materials/apollo-demo-day-transcript.pdf`，对照 Apollo Catalog、release channel、monitor 和 vulnerability 示例；再阅读 Palantir Rubix 当前官方文档，比较课堂 40--72 小时口述与当前 48 小时 node-lifetime 表述。最后选择一个非国防领域的高风险 workflow，例如医院床位、能源调度或工业维护，用 OODA 与 privacy-security frontier 检查其权限和反馈。

\end{document}
